\documentclass[11pt]{article}
\usepackage[margin=1in]{geometry}
\usepackage{amsmath,amssymb,amsthm,booktabs}
\usepackage[hidelinks]{hyperref}
\usepackage{xurl}
\setlength{\emergencystretch}{3em}
\newtheorem{lemma}{Lemma}
\newtheorem{theorem}[lemma]{Theorem}
\newtheorem{corollary}[lemma]{Corollary}
\newcommand{\C}{\mathbb C}
\newcommand{\rr}{r_{\rm refl}}
\newcommand{\rk}{\operatorname{rank}}
\title{A finite complex reflection boundary for sixteen knots}
\author{Written class reduction and primary classification import}
\date{17 November 2026}
\begin{document}
\maketitle

No finite irreducible complex reflection group of rank at most three
can require four generators from a single generating conjugacy class
of reflections. Consequently finite complex reflection-class quotients
cannot certify meridional rank four for the sixteen alternating knots
covered by the retained signed generic obstruction. This is a classification consequence with an
elementary reduction of the class condition. It is not a conclusion
from a failed finite search.

For a finite complex reflection group $W$, let $\rr(W)$ be the minimum
number of reflections generating $W$. If a conjugacy class $C$ of
reflections generates $W$, let $r_C(W)$ be the minimum number of
members of that class generating $W$. The rank is the dimension of the
essential reflection representation.

\begin{lemma}\label{lem:class}
If $W$ is generated by one reflection class $C$, then its reflecting
hyperplanes form one $W$-orbit. Every member of $C$ generates the full
cyclic pointwise stabilizer of its hyperplane, and
\[
                         r_C(W)=\rr(W).
\]
\end{lemma}
\begin{proof}
For an orbit $\mathcal O$ of reflecting hyperplanes, choose nonzero
linear equations $\alpha_H$ and set
$P_{\mathcal O}=\prod_{H\in\mathcal O}\alpha_H$.
Because $W$ permutes these hyperplanes, there is a character
$\chi_{\mathcal O}:W\to\C^*$ satisfying
$P_{\mathcal O}(gv)=\chi_{\mathcal O}(g)P_{\mathcal O}(v)$.

Let $s$ be a reflection with hyperplane $K$ and nontrivial eigenvalue
$\lambda$. If $K\notin\mathcal O$, evaluate at a point of $K$ avoiding
the other hyperplanes to obtain $\chi_{\mathcal O}(s)=1$.
If $K\in\mathcal O$, choose a point $v_0\in K$ on none of the other
hyperplanes and a normal eigenvector $n$ for $s$.
The nonzero linear term of $P_{\mathcal O}(v_0+tn)$ changes by
$\lambda$ under $s$, so $\chi_{\mathcal O}(s)=\lambda\ne1$.
The dimension-one case has only one hyperplane and needs no generic-point
argument.

The hyperplanes of the members of $C$ lie in one orbit.
Any different orbit would therefore give a nontrivial character trivial
on every member of $C$, contradicting generation of $W$.

For a hyperplane $H$, its pointwise stabilizer $W_H$ is finite cyclic:
after averaging a Hermitian form it acts faithfully on the normal line.
Write $|W_H|=e$, choose a primitive generator with determinant
$\zeta_e$, and let the eigenvalue of a member of $C$ on $H$ be
$\zeta_e^k$. Such a member exists since the hyperplanes form one orbit.
The determinant image of $W$ is generated by the common determinant
of $C$. It contains $\zeta_e$, so $\gcd(k,e)=1$.
Hence this member generates $W_H$.

Every reflection with hyperplane $H$ is a power of that member.
Thus an arbitrary reflection generating set $t_1,\ldots,t_r$ can be
replaced by $c_1,\ldots,c_r\in C$, with $t_i\in\langle c_i\rangle$.
The new set still generates $W$. This gives $r_C(W)\le\rr(W)$;
the reverse inequality is immediate.
\end{proof}

\begin{theorem}\label{thm:lowrank}
If $W$ is an irreducible finite complex reflection group of rank at
most three and a reflection class $C$ generates $W$, then $r_C(W)\le3$.
\end{theorem}
\begin{proof}
The original classification gives cyclic rank one, imprimitive
$G(m,p,n)$, symmetric types, and the numbered exceptional groups
\cite[Sections~2--4 and Table~VII]{st54}.
Cyclic rank one needs one generator. The symmetric types $A_2$ and
$A_3$ are generated by two and three adjacent transpositions, respectively.
The primitive rank-two groups need at most three reflections
\cite[printed p.~280]{st54}.

In an imprimitive group with $p<m$, both coordinate and
off-coordinate reflecting hyperplanes occur. Monomial transformations
preserve this distinction, so Lemma~\ref{lem:class} excludes a
generating class. This removes, in particular, the rank-three groups
$G(m,p,3)$ with $1<p<m$ that require four reflections.
It would be false to assert a bound of three for all rank-three
reflection groups without this exclusion.

For $p=m$, the rank-two group is dihedral, generated by two
reflections. In rank three, $G(m,m,3)$, $m\ge2$, is generated by
\[
                       s_{12}(0),\quad s_{23}(0),\quad s_{12}(1),
\]
where $s_{ij}(a)$ interchanges coordinates $i,j$ with reciprocal
$\zeta_m^a,\zeta_m^{-a}$ factors.
The first two generate all permutations. The product
$s_{12}(1)s_{12}(0)$, with permutation conjugates, generates every
diagonal matrix whose exponent sum is zero modulo $m$.
Thus these are three reflection generators.

The rank-three primitive exceptional groups are $G_{23},\ldots,G_{27}$;
each is generated by three reflections. The symmetric type $A_3$
is likewise three-generated. These assertions are explicit in the
rank and generator exception information in Table~VII and its
footnote, printed p.~301 \cite{st54}.
In particular $G_{26}$ is well generated.
Apply Lemma~\ref{lem:class} to each remaining case.
\end{proof}

\begin{corollary}\label{cor:knots}
Let $K$ be any of
\[
\begin{aligned}
\mathcal K_{16}=\{&
13a_{373},13a_{375},13a_{391},13a_{392},\\
&13a_{3264},13a_{3266},13a_{3275},13a_{3276},\\
&13a_{304},13a_{362},13a_{623},13a_{629},\\
&13a_{3232},13a_{3259},13a_{3374},13a_{3375}\}.
\end{aligned}
\]
For every finite-image complex representation of its knot group whose
meridian has displacement rank one, the meridian conjugacy class in
the image generates that image with at most three members.
The ambient dimension is unrestricted.
\end{corollary}
\begin{proof}
The knot group is normally generated by a meridian.
A finite complex image is completely reducible, as can be proved by
averaging any projection over that finite image.
In an irreducible decomposition, the ranks of the meridian
displacements add to one. Thus exactly one irreducible summand $U$
is nontrivial on the meridian. Normal generation makes the whole image
trivial on every other summand.
The restriction to $U$ is therefore faithful on the image.

The imported statement of \path{GENERIC_KCH_EXTENSIONS.tex},
Corollary~\texttt{cor:semisimple}, is the following: for each of these
sixteen knots, every finite-dimensional completely reducible complex
representation $\rho$ with $\rk(\rho(m)-I)=1$ satisfies
$\dim V-\dim V^{\pi_1(S^3\setminus K)}\le3$.
Here that codimension equals $\dim U$, so $\dim U\le3$.
The image on $U$ is an irreducible finite complex reflection group.
Normal generation says its meridian conjugacy class generates it.
Theorem~\ref{thm:lowrank} applies.
\end{proof}

Thus a sharp-four certificate on these knots cannot come from a finite
quotient whose meridian class has a faithful complex-reflection model.
This statement concerns complex representations. It excludes no modular
or nonsemisimple argument, and no finite quotient lacking such a
reflection model. In particular it does not exclude the retained affine
finite-group methods with larger meridian displacement rank.

\section*{Primary import and assurance}

The retained original primary facsimile is
\path{sources/shephard-todd-finite-unitary-1954.pdf}.
Its classification is imported at the precise scope used above.
The complete list, the rank-three rows and the exceptional-generator
footnote in Table~VII were inspected directly, as was the rank-two
statement on printed p.~280. The audit and exact source hash are retained
in \path{finite_complex_reflection_source_audit.json}.
The original classification is not reproved by this note.

The orbit-character and power-replacement argument is written in full.
It was independently challenged by another internal agent.
The knot corollary additionally imports only the stated semisimple
consequence of the signed generic obstruction.
These are written proofs with primary dependencies, not conclusions
from parameter sampling. No formal verification, external specialist
review or worldwide novelty claim is made.

\begin{thebibliography}{9}
\bibitem{st54}
G. C. Shephard and J. A. Todd,
Finite unitary reflection groups,
Canad. J. Math. 6 (1954), 274--304.
Sections~2--4, especially pp.~276--280;
Table~VII and footnote, p.~301.
\url{https://doi.org/10.4153/CJM-1954-028-3}.
Retained original facsimile:
\url{https://math.ucsd.edu/~nwallach/shephard-todd.pdf}.
\end{thebibliography}
\end{document}
